% claim.tex -- AI4Math claim of record (AI-generated, 2026-09-28; instantiated
% from harness/templates/claim/claim.tex per SOP 08b).
% Machine-readable theorem anchors are the "% LEAN:" lines; scripts/paper-lint.sh
% and the claim audit check them. All Lean listings are verbatim, byte-identical
% contiguous excerpts of frozen artifacts of tasks/20260928-comb09-4ntwocorner:
%   statements : formalized/01-comb09-4ntwocorner-statements.lean (127 lines, LF,
%                sha256 4cb1083846bcc64e249717309e63d071d0ecf5bb03f23d85affeaab8152ce005)
%   final      : final/01-comb09-4ntwocorner-proved.lean (420 lines, CRLF,
%                sha256 f1aa7099697fe5706bc5f4dd7124d06f4f92d6c2273b349b862c4b963806f53f)
% Line-ending note: the final artifact is stored with CRLF line endings (the fact
% is disclosed in the task card of the source repository, not shipped with this
% package); the proof listings below normalize line endings to LF and change no
% other byte. The eight frozen declarations were re-checked byte-identical
% between the two files (modulo line endings) by the paper worker on 2026-09-28.
% Line ranges used: statement 30-47, 49-87, 89-93, 96-107, 110-126; final 106-150,
% 155-194, 200-228, 283-311 (the same nine blocks verified in the paper package).
% Verification: audit/listing-verify.txt (audit/inspect-tex.py --verify).
\documentclass[11pt]{article}

\usepackage[T1]{fontenc}
\usepackage[utf8]{inputenc}
\usepackage{amsmath,amssymb,amsthm}
\usepackage{graphicx}
\usepackage{hyperref}
\usepackage{xurl}
% Typography only: xurl lets the long artifact paths and URLs in prose break
% anywhere; emergencystretch absorbs the rest. No listing byte is affected.
\setlength{\emergencystretch}{3em}
% lstlean.tex — Lean style for the listings package.
% Lineage: Jeremy Avigad (2015), adapted from Assia Mahboubi's SSR style;
% inspired by Olivier Verdier's unixode.sty for the Unicode literate table.
% No CTAN package or upstream repo exists; papers carry this file in their
% source package (cap set arXiv:1907.01449, sphere eversion arXiv:2210.07746).
% AI4Math adaptation (AI-generated, 2026-09-26): sorry/admit/sorryAx elevated
% to a dedicated error tier, matching the pipeline's 0-sorry constitution.
\usepackage{listings}
\usepackage{xcolor}

\definecolor{keywordcolor}{rgb}{0.7, 0.1, 0.1}   % red keywords
\definecolor{tacticcolor}{rgb}{0.0, 0.1, 0.6}    % blue tactics
\definecolor{commentcolor}{rgb}{0.4, 0.4, 0.4}   % grey comments
\definecolor{symbolcolor}{rgb}{0.0, 0.1, 0.6}    % blue symbols
\definecolor{sortcolor}{rgb}{0.1, 0.5, 0.1}      % green sorts
\definecolor{errorcolor}{rgb}{0.6, 0, 0}         % dark red: sorry/admit
\definecolor{stringcolor}{rgb}{0.5, 0.1, 0.5}    % purple strings

\lstdefinelanguage{lean}{
  % Anything between $ becomes LaTeX math mode
  mathescape=false,
  % Comments may or not include Latex commands
  texcl=false,
  % keywords, list taken from lean-syntax.txt
  morekeywords=[1]{import, prelude, protected, private, noncomputable,
    definition, meta, renaming, hiding, parameter, parameters, begin,
    constant, constants, lemma, variable, variables, theory, print,
    theorem, example, open, section, namespace, universe, universes, end,
    modifier, modifiers, using, export, exposing, axiom, inductive,
    coinductive, with, structure, class, instance, attribute, local, where,
    by, have, show, let, in, if, then, else, match, do, for, fun, assume,
    from, take, calc, include, omit, infix, infixl, infixr, notation,
    macro, syntax, elab, set_option, initialize, deriving, opaque, abbrev,
    mutual, partial, scoped, termination_by, decreasing_by},
  % Sorts
  morekeywords=[2]{Sort, Type, Prop},
  % tactics, list taken from lean-syntax.txt (tactic keywords)
  morekeywords=[3]{intro, intros, apply, exact, refine, rfl, simp, simp_all,
    simpa, simp_rw, rw, rewrite, erewrite, ring, ring_nf, omega, decide,
    norm_num, norm_cast, induction, cases, rcases, obtain, constructor,
    ext, funext, conv, congr, field_simp, linarith, nlinarith, positivity,
    tauto, trivial, aesop, use, by_cases, by_contra, push_neg, contrapose,
    symm, trans, assumption, contradiction, exfalso, specialize,
    generalize, revert, clear, rename_i, subst, unfold, delta, change,
    convert, split_ifs, interval_cases, fin_cases, zify, gcongr,
    nth_rewrite, suffices, generalizing, at, only, native_decide},
  % warnings - constitution tier: must never survive into a final proof
  morekeywords=[4]{sorry, admit, sorryAx},
  literate=
    {α}{{\ensuremath{\alpha}}}1
    {β}{{\ensuremath{\beta}}}1
    {γ}{{\ensuremath{\gamma}}}1
    {δ}{{\ensuremath{\delta}}}1
    {ε}{{\ensuremath{\varepsilon}}}1
    {ζ}{{\ensuremath{\zeta}}}1
    {η}{{\ensuremath{\eta}}}1
    {θ}{{\ensuremath{\theta}}}1
    {ι}{{\ensuremath{\iota}}}1
    {κ}{{\ensuremath{\kappa}}}1
    {λ}{{\ensuremath{\lambda}}}1
    {μ}{{\ensuremath{\mu}}}1
    {ν}{{\ensuremath{\nu}}}1
    {ξ}{{\ensuremath{\xi}}}1
    {π}{{\ensuremath{\pi}}}1
    {ρ}{{\ensuremath{\rho}}}1
    {σ}{{\ensuremath{\sigma}}}1
    {τ}{{\ensuremath{\tau}}}1
    {υ}{{\ensuremath{\upsilon}}}1
    {φ}{{\ensuremath{\varphi}}}1
    {χ}{{\ensuremath{\chi}}}1
    {ψ}{{\ensuremath{\psi}}}1
    {ω}{{\ensuremath{\omega}}}1
    {Γ}{{\ensuremath{\Gamma}}}1
    {Δ}{{\ensuremath{\Delta}}}1
    {Θ}{{\ensuremath{\Theta}}}1
    {Λ}{{\ensuremath{\Lambda}}}1
    {Ξ}{{\ensuremath{\Xi}}}1
    {Π}{{\ensuremath{\Pi}}}1
    {Σ}{{\ensuremath{\Sigma}}}1
    {Φ}{{\ensuremath{\Phi}}}1
    {Ψ}{{\ensuremath{\Psi}}}1
    {Ω}{{\ensuremath{\Omega}}}1
    {ℵ}{{\ensuremath{\aleph}}}1
    {∀}{{\ensuremath{\forall}}}1
    {∃!}{{\ensuremath{\exists}!}}1
    {∃}{{\ensuremath{\exists}}}1
    {∈}{{\ensuremath{\in}}}1
    {∉}{{\ensuremath{\notin}}}1
    {∘}{{\ensuremath{\circ}}}1
    {∩}{{\ensuremath{\cap}}}1
    {∪}{{\ensuremath{\cup}}}1
    {⊆}{{\ensuremath{\subseteq}}}1
    {⊂}{{\ensuremath{\subset}}}1
    {⊇}{{\ensuremath{\supseteq}}}1
    {⊃}{{\ensuremath{\supset}}}1
    {⊄}{{\ensuremath{\nsubseteq}}}1
    {∅}{{\ensuremath{\emptyset}}}1
    {∧}{{\ensuremath{\wedge}}}1
    {∨}{{\ensuremath{\vee}}}1
    {¬}{{\ensuremath{\neg}}}1
    {⊤}{{\ensuremath{\top}}}1
    {⊥}{{\ensuremath{\bot}}}1
    {↔}{{\ensuremath{\leftrightarrow}}}1
    {→}{{\ensuremath{\rightarrow}}}1
    {←}{{\ensuremath{\leftarrow}}}1
    {↑}{{\ensuremath{\uparrow}}}1
    {⇒}{{\ensuremath{\Rightarrow}}}1
    {⇐}{{\ensuremath{\Leftarrow}}}1
    {⟹}{{\ensuremath{\Longrightarrow}}}1
    {↦}{{\ensuremath{\mapsto}}}1
    {≤}{{\ensuremath{\leq}}}1
    {≥}{{\ensuremath{\geq}}}1
    {≠}{{\ensuremath{\neq}}}1
    {≈}{{\ensuremath{\approx}}}1
    {≡}{{\ensuremath{\equiv}}}1
    {≃}{{\ensuremath{\simeq}}}1
    {≅}{{\ensuremath{\cong}}}1
    {×}{{\ensuremath{\times}}}1
    {·}{{\ensuremath{\cdot}}}1
    {±}{{\ensuremath{\pm}}}1
    {⊕}{{\ensuremath{\oplus}}}1
    {⊗}{{\ensuremath{\otimes}}}1
    {⋆}{{\ensuremath{\star}}}1
    {▸}{{\ensuremath{\blacktriangleright}}}1
    {⟨}{{\ensuremath{\langle}}}1
    {⟩}{{\ensuremath{\rangle}}}1
    {⦃}{{\ensuremath{\{\!\mid}}}1
    {⦄}{{\ensuremath{\mid\!\}}}}1
    {⟦}{{\ensuremath{[\![}}}1
    {⟧}{{\ensuremath{]\!]}}}1
    {⌊}{{\ensuremath{\lfloor}}}1
    {⌋}{{\ensuremath{\rfloor}}}1
    {⌈}{{\ensuremath{\lceil}}}1
    {⌉}{{\ensuremath{\rceil}}}1
    {√}{{\ensuremath{\surd}}}1
    {∣}{{\ensuremath{\mid}}}1
    {∤}{{\ensuremath{\nmid}}}1
    {∎}{{\ensuremath{\blacksquare}}}1
    {⋯}{{\ensuremath{\cdots}}}1
    {…}{{\ensuremath{\ldots}}}1
    {∑}{{\ensuremath{\sum}}}1
    {∏}{{\ensuremath{\prod}}}1
    {⁻¹}{{\ensuremath{^{-1}}}}1
    {¹}{{\ensuremath{^1}}}1
    {²}{{\ensuremath{^2}}}1
    {³}{{\ensuremath{^3}}}1
    {ⁿ}{{\ensuremath{^n}}}1
    {ᵐ}{{\ensuremath{^m}}}1
    {₀}{{\ensuremath{_0}}}1
    {₁}{{\ensuremath{_1}}}1
    {₂}{{\ensuremath{_2}}}1
    {₃}{{\ensuremath{_3}}}1
    {₄}{{\ensuremath{_4}}}1
    {₅}{{\ensuremath{_5}}}1
    {ₙ}{{\ensuremath{_n}}}1
    {ₘ}{{\ensuremath{_m}}}1
    {ₖ}{{\ensuremath{_k}}}1
    {ᵢ}{{\ensuremath{_i}}}1
    {ⱼ}{{\ensuremath{_j}}}1
    {ℤ}{{\ensuremath{\mathbb{Z}}}}1
    {ℕ}{{\ensuremath{\mathbb{N}}}}1
    {ℚ}{{\ensuremath{\mathbb{Q}}}}1
    {ℝ}{{\ensuremath{\mathbb{R}}}}1
    {ℂ}{{\ensuremath{\mathbb{C}}}}1
    {𝔽}{{\ensuremath{\mathbb{F}}}}1
    {𝒫}{{\ensuremath{\mathcal{P}}}}1
    {∗}{{\ensuremath{\ast}}}1
    {•}{{\ensuremath{\bullet}}}1,
  % Comments
  morecomment=[l]{--},
  morecomment=[s]{/-}{-/},
  morestring=[b]",
  stringstyle=\color{stringcolor},
  keepspaces=true,
  keywordstyle=[1]\color{keywordcolor}\bfseries,
  keywordstyle=[2]\color{sortcolor}\bfseries,
  keywordstyle=[3]\color{tacticcolor}\bfseries,
  keywordstyle=[4]\color{errorcolor}\bfseries\underline,
  escapeinside={(*@}{@*)},
  columns=fullflexible,
  basicstyle=\ttfamily\small,
  commentstyle=\color{commentcolor}\itshape,
  breaklines=true,
  showstringspaces=false,
  frame=single,
  framesep=2mm,
  xleftmargin=4mm,
  numbers=left,
  numberstyle=\tiny\color{commentcolor},
  numbersep=6pt,
}

% Inline Lean code in prose: \lean{Int.ModEq.pow}
\newcommand{\lean}[1]{\lstinline[language=lean,basicstyle=\ttfamily\normalsize]|#1|}

\lstset{language=lean}

% --- XeTeX rendering patch (template lineage papers/cyl3-mod23; extended for
% this note's frozen sources and its verbatim C9/fence quote blocks) ---
% tectonic/XeTeX does not apply the listings literate table to multi-byte
% characters (missing-glyph warnings observed in this pipeline on 2026-09-26/27),
% so the Unicode set used by the listings of THIS note is mapped to math glyphs
% as active characters through the standard \lccode/\lowercase idiom: ASCII-only
% source, no extra packages. The set is the union of the template block and the
% code points that occur in the frozen sources but are absent from the template
% patch (00B2 21A6 2212 27FA), plus 00D7, which occurs in the verbatim
% occupancy-fence quote block. CJK prose inside docstrings and quote blocks is
% routed through xeCJK + SimSun, as in the previous batches. Rendering only:
% listing source bytes are untouched (frozen artifacts), so the byte-exactness
% checks are unaffected. The \ifdefined guard leaves arXiv's pdflatex path
% identical to the template.
% NOTE (paper package, first compile round): U+2026 (horizontal ellipsis) is NOT
% in this patch --- xeCJK lists it in its LongPunct table and aborts at load time
% with "Missing \endcsname inserted" if the character is already active (same
% failure class as 00B7/2014 in the 2026-09-27 batch). It renders through the
% CJK font instead.
\ifdefined\XeTeXversion
  {\lccode`\~="2022 \lowercase{\global\catcode`~=\active \global\def~{\ensuremath{\bullet}}}}
  {\lccode`\~="2115 \lowercase{\global\catcode`~=\active \global\def~{\ensuremath{\mathbb{N}}}}}
  {\lccode`\~="2124 \lowercase{\global\catcode`~=\active \global\def~{\ensuremath{\mathbb{Z}}}}}
  {\lccode`\~="2190 \lowercase{\global\catcode`~=\active \global\def~{\ensuremath{\leftarrow}}}}
  {\lccode`\~="2192 \lowercase{\global\catcode`~=\active \global\def~{\ensuremath{\rightarrow}}}}
  {\lccode`\~="2194 \lowercase{\global\catcode`~=\active \global\def~{\ensuremath{\leftrightarrow}}}}
  {\lccode`\~="2200 \lowercase{\global\catcode`~=\active \global\def~{\ensuremath{\forall}}}}
  {\lccode`\~="2203 \lowercase{\global\catcode`~=\active \global\def~{\ensuremath{\exists}}}}
  {\lccode`\~="2227 \lowercase{\global\catcode`~=\active \global\def~{\ensuremath{\wedge}}}}
  {\lccode`\~="2228 \lowercase{\global\catcode`~=\active \global\def~{\ensuremath{\vee}}}}
  {\lccode`\~="2261 \lowercase{\global\catcode`~=\active \global\def~{\ensuremath{\equiv}}}}
  {\lccode`\~="2264 \lowercase{\global\catcode`~=\active \global\def~{\ensuremath{\leq}}}}
  {\lccode`\~="2265 \lowercase{\global\catcode`~=\active \global\def~{\ensuremath{\geq}}}}
  {\lccode`\~="22A2 \lowercase{\global\catcode`~=\active \global\def~{\ensuremath{\vdash}}}}
  {\lccode`\~="27E8 \lowercase{\global\catcode`~=\active \global\def~{\ensuremath{\langle}}}}
  {\lccode`\~="27E9 \lowercase{\global\catcode`~=\active \global\def~{\ensuremath{\rangle}}}}
  % code points present in this note's frozen sources or verbatim quote blocks,
  % missing from the template patch. Restricted to symbols outside xeCJK's own
  % punctuation table: 00B7, 2014, 2026 and the CJK/fullwidth punctuation also
  % occur but belong to the xeCJK punctuation table, so they are rendered
  % through the CJK font instead (same arrangement as the 2026-09-27 batch).
  {\lccode`\~="00B2 \lowercase{\global\catcode`~=\active \global\def~{\ensuremath{{}^2}}}}
  {\lccode`\~="00D7 \lowercase{\global\catcode`~=\active \global\def~{\ensuremath{\times}}}}
  {\lccode`\~="21A6 \lowercase{\global\catcode`~=\active \global\def~{\ensuremath{\mapsto}}}}
  {\lccode`\~="2212 \lowercase{\global\catcode`~=\active \global\def~{\ensuremath{-}}}}
  {\lccode`\~="27FA \lowercase{\global\catcode`~=\active \global\def~{\ensuremath{\Longleftrightarrow}}}}
\fi
\ifdefined\XeTeXversion
  \usepackage{xeCJK}
  \setCJKmainfont{SimSun}
\fi

\newtheorem{theorem}{Theorem}
\theoremstyle{definition}
\newtheorem{conjecture}{Conjecture}

\title{Matching-count sequence of the $4\times n$ grid graph with two corner
vertices removed: interleave identity and parity/mod-$4$ laws\\ (Lean 4
machine-checked claim of record)}
\author{Aurora0134\thanks{Contact: via the GitHub account Aurora0134; ORCID
placeholder --- to be filled in by the author before any upload. No personal
information is carried in this note.}}
\date{September 28, 2026}

\begin{document}
\maketitle

\begin{abstract}
Let $a(n)$ be the number of matchings (the empty one included) of the graph
obtained from the $4\times n$ grid by deleting the upper and the lower corner
vertex of the last column; exact enumeration starts $2, 15, 209, 2426, 29566,
\dots$, and the sequence satisfies an order-$9$ integer recurrence whose two
position-parity subsequences satisfy the \emph{same} order-$9$ recurrence with
their own initial values. We state and machine-check, in Lean 4 (v4.34.0,
mathlib v4.34.0), three theorems about these recurrence-defined sequences: the
main sequence is exactly the interleaving of the two subsequences; its parity
depends only on the index modulo $5$ (odd if and only if the index is $2$ or
$3$ modulo $5$); and its residue modulo $4$ depends only on the index modulo
$10$, with the ten-term pattern $(2,3,1,2,2,0,1,1,0,0)$. Every proof compiles
with no \texttt{sorry} and the axiom audit reports only \{propext, Quot.sound,
Classical.choice\}. Novelty adjudication: no occupying record found in three
consistent rounds of the pipeline's exclusivity review; value tier \emph{new
sequence} only --- the order-$9$ recurrence is the namesake recurrence of OEIS
A033507 and the perfect-matching subcolumn is OEIS A129113, so no novelty is
claimed for either. The combinatorial identification is not kernel-proved and is
stated as a conjecture. This note is a priority claim of record, not a full
paper.
\end{abstract}

\section{Introduction}

Three theorems about three integer sequences defined by initial values plus
linear recurrences are reported here; the sequences come from counting matchings
of the $4\times n$ grid graph with the upper and lower corner vertices of the
last column removed, and the order-$9$ recurrence was obtained by exact rational
minimal-order fitting of enumerated data (the Berlekamp--Massey family of
methods \cite{massey}); the matching-count context goes back to Heilmann and
Lieb and to Godsil and Gutman \cite{heilmann-lieb,godsil-gutman}. The proofs
were produced by the AI4Math pipeline, in which LLM workers generate statements
and proof scripts and the Lean 4 kernel is the sole adjudicator at every gate:
statements are frozen by hash before proving, and a separate audit re-computes
the symmetric difference, scans the whole file for \texttt{sorry}/\texttt{admit}
including comments, and re-prints the axiom dependencies \cite{lean4,mathlib}.
This note is a priority claim of record deposited so that a dated, citable
public record exists ahead of any narrative paper; the fuller paper for the same
three theorems is prepared separately in the source repository under
\texttt{papers/grid4n-col-2notch/} and the two lanes are linked through the
deposit DOI. No literature review is attempted here.

\section{Statement}\label{sec:statement}

In words, the counted object is the $4$-by-$n$ grid of squares with the top and
bottom corner vertices of the last column removed: one counts the ways to choose
some of the edges so that no two chosen edges share an endpoint, with choosing
none counted as one way. The count begins $2, 15, 209, 2426, 29566, \dots$.
That this sequence is the matching count of that graph is supported by exact
enumeration and is stated here only as Conjecture~\ref{conj:bridge}; the frozen
statement layer contains no graph, no grid and no matching.

Throughout, the three sequences are \emph{defined} by the displayed initial
values and recurrences. The formal layer is $0$-indexed, with the index
convention carried in the frozen definitions' header comments:
$\texttt{atc}(n) = a(n+1)$ for the main sequence, $\texttt{ot}(k) = a(2k+1)$
for the odd-position subsequence and $\texttt{et}(k) = a(2k+2)$ for the
even-position subsequence. The main sequence satisfies the order-$9$ recurrence
\begin{align}
  a(n) = 9a(n-1) &+ 41a(n-2) - 41a(n-3) - 111a(n-4) + 91a(n-5)
  \nonumber\\&\quad + 29a(n-6) - 23a(n-7) - a(n-8) + a(n-9)
  \qquad (n \ge 10),
  \label{eq:recmain}
\end{align}
with initial values $a(1..9) = 2$, $15$, $209$, $2426$, $29566$, $355504$,
$4290501$, $51728089$, $623832332$; each position subsequence satisfies the
\emph{same} order-$9$ recurrence
\begin{equation}\label{eq:recsub}
  \begin{aligned}
    b(k) = {}& 163b(k-1) - 2641b(k-2) + 12479b(k-3) - 22577b(k-4)\\
      & + 16705b(k-5) - 5331b(k-6) + 769b(k-7) - 47b(k-8) + b(k-9)
  \end{aligned}
\end{equation}
for $b \in \{o, e\}$, with initial values the odd- and even-positioned terms of
$a$. Structurally, the subsequence recurrence is what the main recurrence
becomes on the square-conjugate face: it is annihilated by $p(x)p(-x)$, where
$p$ is the characteristic polynomial of \eqref{eq:recmain}, which has the same
order $9$ --- in contrast to order-dropping interleavings, where the subsequence
recurrence would have strictly smaller order. That the subsequence recurrence
used here is exactly this pullback is the content of Theorem~\ref{thm:t1}.

Listings~\ref{lst:atcdef} and~\ref{lst:otetdef} show the three sequence
definitions verbatim.

\begin{lstlisting}[caption={The main sequence \lean{atc}: nine initial values and
the order-9 recurrence of \eqref{eq:recmain}, verbatim from the frozen statement
file \texttt{formalized/01-comb09-4ntwocorner-statements.lean} (lines 30--47,
sha256 \texttt{4cb10838}). The docstring carries the $0$-based convention
$\texttt{atc}(n) = a(n+1)$.},label={lst:atcdef}]
/-- **序列 atc（主列，九阶全滞后递推）**。
初值 a(0..8) = 2, 15, 209, 2426, 29566, 355504, 4290501, 51728089, 623832332
（0-based；1-based 即 a(1..9)，逐位取自卡面与 counts.txt）；
递推 atc(n+9) = 9·atc(n+8) + 41·atc(n+7) − 41·atc(n+6) − 111·atc(n+5)
  + 91·atc(n+4) + 29·atc(n+3) − 23·atc(n+2) − atc(n+1) + atc(n)（n ≥ 0）。 -/
def atc : ℕ → ℤ
  | 0 => 2
  | 1 => 15
  | 2 => 209
  | 3 => 2426
  | 4 => 29566
  | 5 => 355504
  | 6 => 4290501
  | 7 => 51728089
  | 8 => 623832332
  | n + 9 =>
      9 * atc (n + 8) + 41 * atc (n + 7) - 41 * atc (n + 6) - 111 * atc (n + 5)
        + 91 * atc (n + 4) + 29 * atc (n + 3) - 23 * atc (n + 2) - atc (n + 1) + atc n
\end{lstlisting}

\begin{lstlisting}[caption={The two position subsequences \lean{ot} (odd side) and
\lean{et} (even side), verbatim from the frozen statement file (lines 49--87,
sha256 \texttt{4cb10838}). Both satisfy the same order-9 recurrence
\eqref{eq:recsub}; their initial values are the odd- and even-positioned terms
of \lean{atc}.},label={lst:otetdef}]
/-- **序列 ot（奇位侧子列，九阶递推）**：ot k = a(2k+1)。
初值 o(0..8) = 2, 209, 29566, 4290501, 623832332, 90717725366,
  13192325177001, 1918451920976894, 278984786074420317
（= atc 在下标 0,2,…,16 处之值，即 1-based a(1),a(3),…,a(17)）；
递推 ot(k+9) = 163·ot(k+8) − 2641·ot(k+7) + 12479·ot(k+6) − 22577·ot(k+5)
  + 16705·ot(k+4) − 5331·ot(k+3) + 769·ot(k+2) − 47·ot(k+1) + ot(k)（k ≥ 0）。 -/
def ot : ℕ → ℤ
  | 0 => 2
  | 1 => 209
  | 2 => 29566
  | 3 => 4290501
  | 4 => 623832332
  | 5 => 90717725366
  | 6 => 13192325177001
  | 7 => 1918451920976894
  | 8 => 278984786074420317
  | k + 9 =>
      163 * ot (k + 8) - 2641 * ot (k + 7) + 12479 * ot (k + 6) - 22577 * ot (k + 5)
        + 16705 * ot (k + 4) - 5331 * ot (k + 3) + 769 * ot (k + 2) - 47 * ot (k + 1) + ot k

/-- **序列 et（偶位侧子列，九阶递推）**：et k = a(2k+2)。
初值 e(0..8) = 15, 2426, 355504, 51728089, 7522727192, 1093972462647,
  159087517952550, 23134798325751388, 3364304696158749633
（= atc 在下标 1,3,…,17 处之值，即 1-based a(2),a(4),…,a(18)）；
递推 et(k+9) = 163·et(k+8) − 2641·et(k+7) + 12479·et(k+6) − 22577·et(k+5)
  + 16705·et(k+4) − 5331·et(k+3) + 769·et(k+2) − 47·et(k+1) + et(k)（k ≥ 0）。 -/
def et : ℕ → ℤ
  | 0 => 15
  | 1 => 2426
  | 2 => 355504
  | 3 => 51728089
  | 4 => 7522727192
  | 5 => 1093972462647
  | 6 => 159087517952550
  | 7 => 23134798325751388
  | 8 => 3364304696158749633
  | k + 9 =>
      163 * et (k + 8) - 2641 * et (k + 7) + 12479 * et (k + 6) - 22577 * et (k + 5)
        + 16705 * et (k + 4) - 5331 * et (k + 3) + 769 * et (k + 2) - 47 * et (k + 1) + et k
\end{lstlisting}

% LEAN: tasks/20260928-comb09-4ntwocorner :: atc_interleave grade=完全证明
\begin{theorem}[interleaving]\label{thm:t1}
For every $k \in \mathbb{N}$, the main sequence splits by position parity:
\[
  \texttt{atc}(2k) = \texttt{ot}(k)
  \qquad\text{and}\qquad
  \texttt{atc}(2k+1) = \texttt{et}(k).
\]
In the $1$-based counting convention: $a(2k+1) = \texttt{ot}(k)$ and
$a(2k+2) = \texttt{et}(k)$.
\end{theorem}

\begin{lstlisting}[caption={Formal statement of Theorem~\ref{thm:t1}, verbatim
from the frozen statement file \texttt{formalized/01-comb09-4ntwocorner-statements.lean}
(lines 89--93, sha256 \texttt{4cb10838}); the \texttt{:= by} opens the proof body,
which is given in Listing~\ref{lst:t1proof}.},label={lst:t1stmt}]
/-- **T1（主攻）交织恒等式**：完整九阶递推 = 两条同签名九阶子列的交织——
对一切 k，atc(2k) = ot k 且 atc(2k+1) = et k
（子列递推 = 主递推特征多项式的平方合取 p(x)p(−x) 面，无阶下降）。 -/
theorem atc_interleave (k : ℕ) :
    atc (2 * k) = ot k ∧ atc (2 * k + 1) = et k := by
\end{lstlisting}

% LEAN: tasks/20260928-comb09-4ntwocorner :: atc_mod2_period5 grade=完全证明
\begin{theorem}[parity]\label{thm:t2}
For every $n, r \in \mathbb{N}$ with $n \bmod 5 = r$,
$\texttt{atc}(n) \bmod 2 = \texttt{pat2}(r)$, where \lean{pat2} is the five-bit
pattern of Listing~\ref{lst:t2blk}. Equivalently, in the $1$-based convention,
$a(n)$ is odd if and only if $n \equiv 2$ or $3 \pmod 5$.
\end{theorem}

% LEAN: tasks/20260928-comb09-4ntwocorner :: atc_mod4_period10 grade=完全证明
\begin{theorem}[residue modulo $4$]\label{thm:t3}
For every $n, r \in \mathbb{N}$ with $n \bmod 10 = r$,
$\texttt{atc}(n) \bmod 4 = \texttt{pat4}(r)$, where \lean{pat4} is the ten-bit
pattern of Listing~\ref{lst:t3blk}: $2, 3, 1, 2, 2, 0, 1, 1, 0, 0$ along the
residue classes $1, \dots, 10$ of the $1$-based convention.
\end{theorem}

\begin{lstlisting}[caption={The parity pattern \lean{pat2} and the formal
statement of Theorem~\ref{thm:t2}, verbatim from the frozen statement file
(lines 96--107, sha256 \texttt{4cb10838}).},label={lst:t2blk}]
/-- **余数样式 pat2**（mod 5 留数类 ↦ mod 2 余数，0-based）：
0↦0, 1↦1, 2↦1，其余↦0（即 atc n 为奇 ⟺ n ≡ 1 或 2 (mod 5)；
1-based 口径：a(k) 为奇 ⟺ k ≡ 2 或 3 (mod 5)）。 -/
def pat2 : ℕ → ℤ
  | 0 => 0
  | 1 => 1
  | 2 => 1
  | _ => 0

/-- **T2（伴随）**：atc 的 mod 2 余数仅依赖 n mod 5，样式为 pat2。 -/
theorem atc_mod2_period5 (n r : ℕ) (hr : n % 5 = r) :
    atc n % 2 = pat2 r := by
\end{lstlisting}

\begin{lstlisting}[caption={The mod-4 pattern \lean{pat4} and the formal
statement of Theorem~\ref{thm:t3}, verbatim from the frozen statement file
(lines 110--126, sha256 \texttt{4cb10838}).},label={lst:t3blk}]
/-- **余数样式 pat4**（mod 10 留数类 ↦ mod 4 余数，0-based）：
r ↦ a(r+1) mod 4，r = 0..9 依次 (2,3,1,2,2,0,1,1,0,0)
（1-based 口径：a(k) mod 4 按 k mod 10，k=1..10 同为 (2,3,1,2,2,0,1,1,0,0)）。 -/
def pat4 : ℕ → ℤ
  | 0 => 2
  | 1 => 3
  | 2 => 1
  | 3 => 2
  | 4 => 2
  | 5 => 0
  | 6 => 1
  | 7 => 1
  | _ => 0

/-- **T3（伴随）**：atc 的 mod 4 余数仅依赖 n mod 10，样式为 pat4。 -/
theorem atc_mod4_period10 (n r : ℕ) (hr : n % 10 = r) :
    atc n % 4 = pat4 r := by
\end{lstlisting}

All three statements are of the pipeline's highest grade, complete proof: no
\texttt{sorry}, axiom dependencies inside the whitelist, and statement
byte-identical to the frozen text. The grade is quoted from the audit record
\texttt{audit/final-comb09-4ntwocorner.txt} of the source task; it is not a
judgement made in this note.

\begin{conjecture}[combinatorial bridge, not proved here]\label{conj:bridge}
For every $n \ge 1$, $\texttt{atc}(n-1)$ equals $a(n)$, the number of matchings
of the graph $G_n$ obtained from the $4 \times n$ grid $P_4 \square P_n$ by
deleting the upper and the lower corner vertex of the last column (so the last
column keeps its two middle vertices).
\end{conjecture}

\section{Proof}\label{sec:proof}

All three proofs are strong induction on the index with the seed window
discharged by computation and the step closed by congruence algebra.
Theorem~\ref{thm:t1} is an equation between recurrences and needs one auxiliary
lemma, \lean{atc_sq}, which exhibits an even-lag recurrence for the main
sequence as an integer linear combination of ten shifted instances of the
defining recurrence; its coefficient list is exactly the subsequence signature
\eqref{eq:recsub}, which is the algebraic face of the $p(x)p(-x)$ structure
noted in Section~\ref{sec:statement}. Listing~\ref{lst:atcsq} gives that lemma
in full and Listing~\ref{lst:t1proof} the complete proof of
Theorem~\ref{thm:t1}. Theorems~\ref{thm:t2} and~\ref{thm:t3} share one
skeleton --- drop the residue parameter into an inner statement, collapse the
recurrence modulo $2$ (respectively $4$), split the residue by \lean{mod_cases}
and close each branch by \lean{decide} --- and are excerpted at the inductive
step with the first residue branch; both carry a raised
\texttt{set\_option maxHeartbeats} at declaration level, because the per-branch
arithmetic accumulates. The complete $420$-line final artifact ships in this
deposit under \texttt{proofs/}.

\begin{lstlisting}[caption={The auxiliary lemma \lean{atc_sq} in full, verbatim
excerpt of the final artifact \texttt{final/01-comb09-4ntwocorner-proved.lean}
(lines 106--150, sha256 \texttt{f1aa7099}; line endings normalized from CRLF, no
other byte changed). The lemma is the only declaration the final file adds to
the frozen statements; it lives in the proof layer and is allowed as such by the
audit.},label={lst:atcsq}]
/-- **辅助引理 atc_sq（平方递推，T1 的关键中间件）**：主列 atc 自身满足
偶滞后 18 的递推（= 子列递推 s(y) 在 y=x² 处的拉回 s(x²) = −p(x)p(−x)，
q(E)a(n) = 0 由主递推的 10 个相邻实例线性组合给出，无需归纳）。
系数证书 (h0..h9) ↦ (1,1,−23,−29,91,111,−41,−41,9,1) 经 python 符号展开
与 a(1..50) 数值逐项核验（r1 准备记录）。 -/
theorem atc_sq (n : ℕ) :
    atc (n + 18) = 163 * atc (n + 16) - 2641 * atc (n + 14) + 12479 * atc (n + 12)
      - 22577 * atc (n + 10) + 16705 * atc (n + 8) - 5331 * atc (n + 6)
        + 769 * atc (n + 4) - 47 * atc (n + 2) + atc n := by
  have e : ∀ m : ℕ, atc (m + 9)
      = 9 * atc (m + 8) + 41 * atc (m + 7) - 41 * atc (m + 6) - 111 * atc (m + 5)
        + 91 * atc (m + 4) + 29 * atc (m + 3) - 23 * atc (m + 2) - atc (m + 1) + atc m :=
    fun m => by simp only [atc]
  have h0 : atc (n + 9) = 9 * atc (n + 8) + 41 * atc (n + 7) - 41 * atc (n + 6)
      - 111 * atc (n + 5) + 91 * atc (n + 4) + 29 * atc (n + 3) - 23 * atc (n + 2)
      - atc (n + 1) + atc n := e n
  have h1 : atc (n + 10) = 9 * atc (n + 9) + 41 * atc (n + 8) - 41 * atc (n + 7)
      - 111 * atc (n + 6) + 91 * atc (n + 5) + 29 * atc (n + 4) - 23 * atc (n + 3)
      - atc (n + 2) + atc (n + 1) := e (n + 1)
  have h2 : atc (n + 11) = 9 * atc (n + 10) + 41 * atc (n + 9) - 41 * atc (n + 8)
      - 111 * atc (n + 7) + 91 * atc (n + 6) + 29 * atc (n + 5) - 23 * atc (n + 4)
      - atc (n + 3) + atc (n + 2) := e (n + 2)
  have h3 : atc (n + 12) = 9 * atc (n + 11) + 41 * atc (n + 10) - 41 * atc (n + 9)
      - 111 * atc (n + 8) + 91 * atc (n + 7) + 29 * atc (n + 6) - 23 * atc (n + 5)
      - atc (n + 4) + atc (n + 3) := e (n + 3)
  have h4 : atc (n + 13) = 9 * atc (n + 12) + 41 * atc (n + 11) - 41 * atc (n + 10)
      - 111 * atc (n + 9) + 91 * atc (n + 8) + 29 * atc (n + 7) - 23 * atc (n + 6)
      - atc (n + 5) + atc (n + 4) := e (n + 4)
  have h5 : atc (n + 14) = 9 * atc (n + 13) + 41 * atc (n + 12) - 41 * atc (n + 11)
      - 111 * atc (n + 10) + 91 * atc (n + 9) + 29 * atc (n + 8) - 23 * atc (n + 7)
      - atc (n + 6) + atc (n + 5) := e (n + 5)
  have h6 : atc (n + 15) = 9 * atc (n + 14) + 41 * atc (n + 13) - 41 * atc (n + 12)
      - 111 * atc (n + 11) + 91 * atc (n + 10) + 29 * atc (n + 9) - 23 * atc (n + 8)
      - atc (n + 7) + atc (n + 6) := e (n + 6)
  have h7 : atc (n + 16) = 9 * atc (n + 15) + 41 * atc (n + 14) - 41 * atc (n + 13)
      - 111 * atc (n + 12) + 91 * atc (n + 11) + 29 * atc (n + 10) - 23 * atc (n + 9)
      - atc (n + 8) + atc (n + 7) := e (n + 7)
  have h8 : atc (n + 17) = 9 * atc (n + 16) + 41 * atc (n + 15) - 41 * atc (n + 14)
      - 111 * atc (n + 13) + 91 * atc (n + 12) + 29 * atc (n + 11) - 23 * atc (n + 10)
      - atc (n + 9) + atc (n + 8) := e (n + 8)
  have h9 : atc (n + 18) = 9 * atc (n + 17) + 41 * atc (n + 16) - 41 * atc (n + 15)
      - 111 * atc (n + 14) + 91 * atc (n + 13) + 29 * atc (n + 12) - 23 * atc (n + 11)
      - atc (n + 10) + atc (n + 9) := e (n + 9)
  linear_combination h9 + 9 * h8 - 41 * h7 - 41 * h6 + 111 * h5 + 91 * h4
    - 29 * h3 - 23 * h2 + h1 + h0
\end{lstlisting}

\begin{lstlisting}[caption={Proof of Theorem~\ref{thm:t1} in full, verbatim
excerpt of the final artifact (lines 155--194, sha256 \texttt{f1aa7099}; line
endings normalized from CRLF).},label={lst:t1proof}]
theorem atc_interleave (k : ℕ) :
    atc (2 * k) = ot k ∧ atc (2 * k + 1) = et k := by
  induction k using Nat.strong_induction_on with
  | h k ih =>
    by_cases hk : k < 9
    · interval_cases k <;> decide
    · obtain ⟨j, rfl⟩ : ∃ j, k = j + 9 := ⟨k - 9, by omega⟩
      constructor
      · have e18 : 2 * (j + 9) = 2 * j + 18 := by ring
        rw [e18, atc_sq]
        simp only [ot]
        rw [show 2 * j + 16 = 2 * (j + 8) by ring,
            show 2 * j + 14 = 2 * (j + 7) by ring,
            show 2 * j + 12 = 2 * (j + 6) by ring,
            show 2 * j + 10 = 2 * (j + 5) by ring,
            show 2 * j + 8 = 2 * (j + 4) by ring,
            show 2 * j + 6 = 2 * (j + 3) by ring,
            show 2 * j + 4 = 2 * (j + 2) by ring,
            show 2 * j + 2 = 2 * (j + 1) by ring,
            (ih (j + 8) (by omega)).1, (ih (j + 7) (by omega)).1,
            (ih (j + 6) (by omega)).1, (ih (j + 5) (by omega)).1,
            (ih (j + 4) (by omega)).1, (ih (j + 3) (by omega)).1,
            (ih (j + 2) (by omega)).1, (ih (j + 1) (by omega)).1,
            (ih j (by omega)).1]
      · have e19 : 2 * (j + 9) + 1 = 2 * j + 1 + 18 := by ring
        rw [e19, atc_sq]
        simp only [et]
        rw [show 2 * j + 1 + 16 = 2 * (j + 8) + 1 by ring,
            show 2 * j + 1 + 14 = 2 * (j + 7) + 1 by ring,
            show 2 * j + 1 + 12 = 2 * (j + 6) + 1 by ring,
            show 2 * j + 1 + 10 = 2 * (j + 5) + 1 by ring,
            show 2 * j + 1 + 8 = 2 * (j + 4) + 1 by ring,
            show 2 * j + 1 + 6 = 2 * (j + 3) + 1 by ring,
            show 2 * j + 1 + 4 = 2 * (j + 2) + 1 by ring,
            show 2 * j + 1 + 2 = 2 * (j + 1) + 1 by ring,
            (ih (j + 8) (by omega)).2, (ih (j + 7) (by omega)).2,
            (ih (j + 6) (by omega)).2, (ih (j + 5) (by omega)).2,
            (ih (j + 4) (by omega)).2, (ih (j + 3) (by omega)).2,
            (ih (j + 2) (by omega)).2, (ih (j + 1) (by omega)).2,
            (ih j (by omega)).2]
\end{lstlisting}

\begin{lstlisting}[caption={Theorem~\ref{thm:t2}: the clean inner lemma, the
induction shell, the mod-2 collapse and the first residue branch, verbatim
excerpt of the final artifact (lines 200--228, sha256 \texttt{f1aa7099}; line
endings normalized from CRLF). The theorem closes at line 277 of the artifact by
rewriting the hypothesis into the inner statement; the other four residue
branches are identical in form with different residues.},label={lst:t2step}]
  have main : ∀ n : ℕ, atc n % 2 = pat2 (n % 5) := by
    intro n
    induction n using Nat.strong_induction_on with
    | h n ih =>
      by_cases hn : n < 9
      · interval_cases n <;> decide
      · obtain ⟨j, rfl⟩ : ∃ j, n = j + 9 := ⟨n - 9, by omega⟩
        have h1 : atc (j + 9) % 2
            = (atc (j + 8) % 2 + atc (j + 7) % 2 - atc (j + 6) % 2 - atc (j + 5) % 2
                + atc (j + 4) % 2 + atc (j + 3) % 2 - atc (j + 2) % 2 - atc (j + 1) % 2
                + atc j % 2) % 2 := by
          simp only [atc]
          omega
        rw [h1, ih (j + 8) (by omega), ih (j + 7) (by omega), ih (j + 6) (by omega),
            ih (j + 5) (by omega), ih (j + 4) (by omega), ih (j + 3) (by omega),
            ih (j + 2) (by omega), ih (j + 1) (by omega), ih j (by omega)]
        mod_cases hj : j % 5
        · have hj0 : j % 5 = 0 := by simpa [Nat.ModEq] using hj
          have hj1 : (j + 1) % 5 = 1 := by omega
          have hj2 : (j + 2) % 5 = 2 := by omega
          have hj3 : (j + 3) % 5 = 3 := by omega
          have hj4 : (j + 4) % 5 = 4 := by omega
          have hj5 : (j + 5) % 5 = 0 := by omega
          have hj6 : (j + 6) % 5 = 1 := by omega
          have hj7 : (j + 7) % 5 = 2 := by omega
          have hj8 : (j + 8) % 5 = 3 := by omega
          have hj9 : (j + 9) % 5 = 4 := by omega
          rw [hj0, hj1, hj2, hj3, hj4, hj5, hj6, hj7, hj8, hj9]
          decide
\end{lstlisting}

\begin{lstlisting}[caption={Theorem~\ref{thm:t3}: the same skeleton at period
$10$ --- the mod-4 collapse and the first residue branch, verbatim excerpt of
the final artifact (lines 283--311, sha256 \texttt{f1aa7099}; line endings
normalized from CRLF).},label={lst:t3step}]
  have main : ∀ n : ℕ, atc n % 4 = pat4 (n % 10) := by
    intro n
    induction n using Nat.strong_induction_on with
    | h n ih =>
      by_cases hn : n < 9
      · interval_cases n <;> decide
      · obtain ⟨j, rfl⟩ : ∃ j, n = j + 9 := ⟨n - 9, by omega⟩
        have h1 : atc (j + 9) % 4
            = (atc (j + 8) % 4 + atc (j + 7) % 4 - atc (j + 6) % 4 + atc (j + 5) % 4
                - atc (j + 4) % 4 + atc (j + 3) % 4 + atc (j + 2) % 4 - atc (j + 1) % 4
                + atc j % 4) % 4 := by
          simp only [atc]
          omega
        rw [h1, ih (j + 8) (by omega), ih (j + 7) (by omega), ih (j + 6) (by omega),
            ih (j + 5) (by omega), ih (j + 4) (by omega), ih (j + 3) (by omega),
            ih (j + 2) (by omega), ih (j + 1) (by omega), ih j (by omega)]
        mod_cases hj : j % 10
        · have hj0 : j % 10 = 0 := by simpa [Nat.ModEq] using hj
          have hj1 : (j + 1) % 10 = 1 := by omega
          have hj2 : (j + 2) % 10 = 2 := by omega
          have hj3 : (j + 3) % 10 = 3 := by omega
          have hj4 : (j + 4) % 10 = 4 := by omega
          have hj5 : (j + 5) % 10 = 5 := by omega
          have hj6 : (j + 6) % 10 = 6 := by omega
          have hj7 : (j + 7) % 10 = 7 := by omega
          have hj8 : (j + 8) % 10 = 8 := by omega
          have hj9 : (j + 9) % 10 = 9 := by omega
          rw [hj0, hj1, hj2, hj3, hj4, hj5, hj6, hj7, hj8, hj9]
          decide
\end{lstlisting}

\section{Verification evidence}

\begin{itemize}
\item Toolchain: Lean \texttt{v4.34.0} (\texttt{leanprover/lean4:v4.34.0}),
mathlib4 \texttt{v4.34.0}, revision \texttt{5ed29652}
\cite{lean4,mathlib,mathlib4-tag}; this note is typeset with tectonic 0.17.0.
\item Statement integrity: the frozen statement file
\texttt{formalized/}\allowbreak\texttt{01-comb09-4ntwocorner-statements.lean}
($127$ lines) has
sha256 \texttt{4cb1083846bcc64e249717309e63d071}\allowbreak
\texttt{d0ecf5bb03f23d85affeaab8152ce005};
the final file \texttt{final/}\allowbreak\texttt{01-comb09-4ntwocorner-proved.lean}
($420$ lines,
CRLF at rest) has sha256
\texttt{f1aa7099697fe5706bc5f4dd7124d06f}\allowbreak
\texttt{4f92d6c2273b349b862c4b963806f53f}. The
final file was compared with the frozen file by symmetric difference over the
declaration region: $8$ of $8$ frozen declarations byte-identical, all $8$
docstrings contained verbatim; the only addition is the proof-layer lemma
\lean{atc_sq}.
\item Kernel check: compiles with $0$ \texttt{sorry} (whole-file scan including
comments, case-insensitive, also for \texttt{admit}/\texttt{sorryAx}); the
strict re-compile exits $0$ with zero diagnostics. The axiom audit output,
verbatim (kept literal; the smaller monospace width is typesetting only):
\begin{lstlisting}[basicstyle=\ttfamily\footnotesize,frame=none]
AXIOMS atc_interleave [propext, Classical.choice, Quot.sound]
AXIOMS atc_mod2_period5 [propext, Quot.sound]
AXIOMS atc_mod4_period10 [propext, Quot.sound]
\end{lstlisting}
\item Grading by the audit department: all three theorems \emph{complete
proof}; lowest grade reached for the case = \emph{complete proof}; no
scaffolding route was used.
\item Audit chain (originals in this deposit under \texttt{audit/}): final
adjudication \texttt{final-comb09-4ntwocorner.txt} (four independent checks),
the axiom printouts \texttt{axioms-atc\_interleave.txt},
\texttt{axioms-atc\_mod2\_period5.txt} and
\texttt{axioms-atc\_mod4\_period10.txt}, the rechecker's independent re-run
\texttt{axioms-recheck.txt}, this note's lint and compile logs and listing
byte-exactness record, and the full C9 search log \texttt{c9-record.md}. The
semantic-fidelity check (roundtrip re-translation) of the underlying task ran on
the same model node as its formalisation; that independence weakness is
mitigated, not removed, by kernel adjudication and manual review.
\end{itemize}

\section{Novelty evidence}\label{sec:novelty}

Adjudication copied from the pipeline's exclusivity review (C9 gate), which ran
before the problem was accepted and was repeated in a third, deepest round by
workers who re-derived the data independently rather than trusting the external
report; all three rounds returned \emph{no occupying record}, and the third
round graded the value tier strictly as \emph{new sequence}, correcting the
earlier ``new sequence and new recurrence'' wording. It is not a claim of new
mathematics and not a first-discovery claim. Summary by channel, from the
third-round record (\texttt{tasks/20260928-comb0709-deeprecheck/recheck/deep-recheck-combm-03.md},
kept in the source repository); the full query log, including every zero-hit
query, is in this deposit at \texttt{audit/c9-record.md}.

\begin{center}
\resizebox{\textwidth}{!}{%
\begin{tabular}{lll}
\hline
channel & queries (round 3) & result \\
\hline
OEIS sequence search \cite{oeis} & 9 new strings & zero hits (the main
sequence accumulated $40+$ independent query strings across rounds: head,
truncations, prepended/zero-padded variants, characteristic-polynomial and
numerator strings, mod-$3$ subsequences, tail windows, difference windows) \\
OEIS per-residue subsequences & included above & zero hits (mod-$3$ residues
newly split in round 3) \\
zbMATH / arXiv / Math StackExchange / Crossref & 20 new alias groups & zero
same-shaped hits; two arXiv probes return \texttt{total=0} \\
literature, closest neighbour \cite{oh2019} & full-text read & Theorem 4 is a
method-general form with a differently-shaped statistic; the sequence does not
occur \\
OEIS record deep read & A033507, A033506 full columns & mother-family
literature covers only undamaged grids; no defect-variant citation \\
library layer & mathlib, compfiles, sequencelib & zero same-shaped objects;
mathlib has the matching predicates \cite{mathlib-matching} but no counting
theorem \\
public Lean ecosystem & code search on newly computed large terms &
\texttt{total=0} \\
\hline
\end{tabular}%
}
\end{center}

The review's machine checks, re-derived independently in round 3: the $12$-term
head was reproduced term by term by three heterogeneous methods; the order-$9$
recurrence holds on $50$ computed terms; the minimal order over $\mathbb{Q}$ is
exactly $9$ (no order-$8$ fit). Two occupancies were established and are carried
as boundaries: A033507 \cite{oeis-a033507}, whose recurrence line is the
namesake recurrence of this note, and A129113 \cite{oeis-a129113} for the
perfect-matching subcolumn.

The C9 verdict of the third round is reproduced verbatim below (the full record
ships as \texttt{audit/c9-record.md} in this deposit):
\begin{lstlisting}[basicstyle=\ttfamily\footnotesize,frame=none]
**C9 四态：查无占位（三轮一致，本审支持前两审结论）。**
**价值分级：新序列**（九阶递推与已挂名的 A033507 完全同谱——前两审已注，本审机器确认；「新序列·新递推」措辞本审维持修正为「新序列」）。
\end{lstlisting}

The occupancy fences of the source claim log
(\texttt{tasks/20260928-}\allowbreak\texttt{comb09-4ntwocorner/}\allowbreak\texttt{phase0/claims.md}
\S 0, kept in the source repository) bound every wording of this note and are
reproduced verbatim below:
\begin{lstlisting}[basicstyle=\ttfamily\footnotesize,frame=none]
**占位围栏：价值级仅「新序列」——九阶递推即 A033507（无缺陷 4×n 全匹配）的挂名递推（同谱姊妹序列），禁止对递推本身作新颖性主张；PM 子列 = OEIS A129113 已挂名占位，禁止以其作新颖性主张。**
\end{lstlisting}

The residual gaps of the review are disclosed with it: the general web search
layer was unreachable in the review environment, so literature coverage is
carried by four structured channels (OpenAlex, zbMATH, arXiv, Crossref); the
Math StackExchange native API was throttled at review time; the mother-family
full texts (Lundow 1996/1998, Read 1982) and Guichard 2008 were not read in
full (the latter only possibly touches the already-occupied perfect-matching
face); and the machine verification of the recurrence stops at $50$ computed
terms. If any of these channels later returns a hit on this sequence, the tier
above weakens accordingly, and this deposit will be amended by a later version
rather than by editing this record.

\section{Scope and limitations}

Grade and boundary, stated as the audit states them. (a) The theorems cover the
recurrence-defined sequences only; their equality with matching counts of the
notched grid is not kernel-proved and is Conjecture~\ref{conj:bridge}. Its
empirical support is enumeration by two independently written programs agreeing
on their overlap; the enumerators themselves are not verified programs, and the
scripts and their outputs stay in the working repository under
\texttt{tasks/20260928-comb09-4ntwocorner/}. (b) Priority boundary: the
order-$9$ recurrence is the namesake recurrence of OEIS A033507
\cite{oeis-a033507} --- the undamaged $4\times n$ grid, a same-spectrum sister
sequence --- so nothing is claimed about the recurrence itself; the
perfect-matching subcolumn is OEIS A129113 \cite{oeis-a129113}, already
occupied, and is claimed by nobody here. The value tier of this note is
strictly ``new sequence'' (initial values) plus the three structural theorems,
and the novelty evidence is negative search evidence with disclosed holes. (c)
Indexing is $0$-based in the formal layer ($\texttt{atc}(n) = a(n+1)$,
$\texttt{ot}(k) = a(2k+1)$, $\texttt{et}(k) = a(2k+2)$); the $1$-based reading
is used for the combinatorial statements. (d) The recurrences are computational
findings: exact rational minimal-order fits (order exactly $9$; no order-$8$
fit) verified against enumeration, not derived from a transfer-matrix argument;
the theorems assert the periods $5$ and $10$ and the pattern tables, not their
minimality, which remains a computed fact. The mod-$3$ period $364$ was
measured the same way and was deliberately not formalised in this batch. (e)
The gate-four sign-off of the underlying task was exercised by the author's
instruction to deposit this claim on 2026-09-28, while the task report file
still carries its draft header (``pending manual sign-off at gate four''); both
facts are recorded here as they stand, and the task's own artifacts are not
edited. (f) The roundtrip semantic check of the underlying task ran on the same
model node as its formalisation; the independence risk is mitigated, not
removed, by kernel adjudication and manual review. (g) This note carries no
literature review and is a deposit record, not a survey; the narrative paper for
the same result is separate.

\section*{Data availability}

The deposit package for this claim contains the claim note (LaTeX source and
PDF), the Lean sources (\texttt{proofs/}: frozen statement file, final proof
file, \texttt{lean-toolchain}), the verification and audit logs (final
adjudication, axiom printouts and recheck, this note's lint and compile logs,
listing byte-exactness record), and the full C9 query log with zero-hit entries
(\texttt{audit/}). Zenodo: the DOI is an unfilled field, \texttt{RESERVED-DOI};
it is reserved by the author at upload and recorded in
\texttt{metadata/README.md} of this package once created; the upload steps are
listed in \texttt{UPLOAD.md} in the source repository, not shipped with this
deposit. Public mirror of this deposit: the claim lane's repository
\url{https://github.com/Aurora0134/ai4math-claims}, directory
\texttt{grid4n-col-2notch/} (pushed under the author's instruction of
2026-09-28; the commit and timestamp are recorded in the source repository's
\texttt{claims/grid4n-col-2notch/card.md}). Claim note under CC BY 4.0, code
under MIT.

\section*{Use of generative AI}

(a) AI performed: candidate generation and the exclusivity review,
autoformalisation of the five definitions and three statements, the roundtrip
semantic check, proof search and proof-script writing, error classification and
repair, the assembly of the final file, the audit re-runs, the prose
translation of the proofs, and the assembly of this note including its verbatim
listings. (b) Model, node and budget: all task-stage sampling ran on one node,
wire-id \texttt{anthropic/a6api-main/kimi-k3}, with no cross-node switching, and
consumed $9$ worker dispatches of a permitted $20$, $1$ LLM call of a permitted
$6$ and $25$ kernel compiles of a permitted $32$, copied verbatim from the task
ledger \texttt{budget.log} (a source repository file, not shipped with this
deposit). This note consumed $0$ pipeline LLM calls; its LaTeX compilations are
counted in this claim lane's own ledger at the same source repository location,
not shipped with this deposit. The node record for 2026-09-28 carries a drift
observation --- the persistent relay record shows both
\texttt{anthropic/a6api-main/kimi-k3} and \texttt{stepfun/step-5-preview} on
that date, and the two independent re-audit dispatches of the paper lane ran on
the latter; this is disclosed rather than smoothed over. (c) Final adjudication:
all statements and proofs reported here were checked by the Lean 4 kernel: every
proof compiles with no \texttt{sorry}, and the axiom audit reports only
\{propext, Quot.sound, Classical.choice\} or a subset thereof for each of
\lean{atc_interleave}, \lean{atc_mod2_period5} and \lean{atc_mod4_period10}.
(d) The AI system is not an author; the human author reviewed the evidence and
is responsible for all content.

\section*{Author contributions}

CRediT-style: the human author adjudicated topic selection, statement semantics
and the deposit decision; the AI4Math pipeline (an LLM) generated the candidate
propositions, the proof searches and the text drafts. The AI system is not
listed as an author.

\begin{thebibliography}{9}
% Reused from papers/grid4n-col-2notch/paper.tex, where each entry was verified
% on 2026-09-28 from the working machine before being typed in: Crossref REST
% metadata for the five DOIs (title and pages matched exactly), the arXiv export
% API plus the landing page for the preprint, the OEIS search endpoint in text
% format for the two sequence entries (their recurrence and data lines read
% directly), and the mathlib4 tag object and raw files for the repository pins.
% Raw responses: papers/grid4n-col-2notch/audit/bib-verify/. No entry is quoted
% from memory; an HTTP 200 on a DOI is not existence evidence (the template's
% Lean 4 DOI ending _27 resolves to a Collatz paper in the same volume; the
% correct ending is _37).
\bibitem{lean4}
L.~de Moura and S.~Ullrich.
\newblock The Lean 4 theorem prover and programming language.
\newblock \emph{Automated Deduction --- CADE-28}, LNCS, pp.~625--635, 2021.
\url{https://doi.org/10.1007/978-3-030-79876-5_37}

\bibitem{mathlib}
The mathlib Community.
\newblock The Lean mathematical library.
\newblock \emph{CPP 2020: Proceedings of the 9th ACM SIGPLAN International
Conference on Certified Programs and Proofs}, pp.~367--381, 2020.
\url{https://doi.org/10.1145/3372885.3373824}

\bibitem{mathlib4-tag}
The mathlib Community.
\newblock The mathlib4 repository, tag \texttt{v4.34.0} (commit
\texttt{5ed2965256430c3649e86755f9576b54eca72435}); its \texttt{lean-toolchain}
pin reads \texttt{leanprover/lean4:v4.34.0}.
\newblock \url{https://github.com/leanprover-community/mathlib4/tree/v4.34.0}

\bibitem{heilmann-lieb}
O.~J. Heilmann and E.~H. Lieb.
\newblock Theory of monomer-dimer systems.
\newblock \emph{Communications in Mathematical Physics} 25 (1972),
pp.~190--232.
\url{https://doi.org/10.1007/BF01877590}

\bibitem{godsil-gutman}
C.~D. Godsil and I.~Gutman.
\newblock On the theory of the matching polynomial.
\newblock \emph{Journal of Graph Theory} 5 (1981), pp.~137--144.
\url{https://doi.org/10.1002/jgt.3190050203}

\bibitem{massey}
J.~L. Massey.
\newblock Shift-register synthesis and BCH decoding.
\newblock \emph{IEEE Transactions on Information Theory} 15 (1969),
pp.~122--127.
\url{https://doi.org/10.1109/TIT.1969.1054260}

\bibitem{oh2019}
S.~Oh.
\newblock State matrix recursion method and monomer--dimer problem.
\newblock \emph{arXiv:1901.07847 [math.CO]}, 2019.
\url{https://arxiv.org/abs/1901.07847}

\bibitem{oeis}
The OEIS Foundation.
\newblock \emph{The On-Line Encyclopedia of Integer Sequences}.
\newblock \url{https://oeis.org/} --- searched for the sequence
\eqref{eq:recmain}'s initial values, its subsequences and variants, with no
entry found (see Section~\ref{sec:novelty}).

\bibitem{oeis-a033507}
The OEIS Foundation.
\newblock Entry A033507 (P.~H. Lundow): number of matchings in graph
$P_4 \times P_n$; recurrence line
$a(n) = 9a(n-1) + 41a(n-2) - 41a(n-3) - 111a(n-4) + 91a(n-5) + 29a(n-6)
- 23a(n-7) - a(n-8) + a(n-9)$ and index-line signature
$(9, 41, -41, -111, 91, 29, -23, -1, 1)$ as displayed.
\newblock \emph{The On-Line Encyclopedia of Integer Sequences}.
\url{https://oeis.org/A033507}

\bibitem{oeis-a129113}
The OEIS Foundation.
\newblock Entry A129113 (R.~L. Bagula, 2007): expansion of
$x^3/(1-x-5x^2-x^3+x^4)$; data
$0, 0, 0, 1, 1, 6, 12, 42, 107, 323, 888, 2568, 7224, 20629, 58429, \dots$
as displayed.
\newblock \emph{The On-Line Encyclopedia of Integer Sequences}.
\url{https://oeis.org/A129113}

\bibitem{mathlib-matching}
The mathlib Community.
\newblock \texttt{Matching.lean}, in \texttt{Mathlib/Combinatorics/}
\newblock \texttt{SimpleGraph/}, mathlib4, tag \texttt{v4.34.0}: defines the
predicates \texttt{IsMatching} and \texttt{IsPerfectMatching} and structural
lemmas, and contains no matching-count function.
\newblock \url{https://github.com/leanprover-community/mathlib4/blob/v4.34.0/Mathlib/Combinatorics/SimpleGraph/Matching.lean}
\end{thebibliography}

\end{document}
